\begin{frame}
\frametitle{Balance primario del GNC}
\centering
\begin{tikzpicture}
\begin{axis}[
    width=12.6cm, height=5.7cm,
    clip=false,
    ylabel={\% del PIB},
    xmin=2019, xmax=2035.25,
    xtick={2019,2021,...,2035},
    xticklabel style={/pgf/number format/.cd,1000 sep={},fixed,precision=0},
    scaled x ticks=false,
    scaled y ticks=false,
    tick label style={font=\footnotesize},
    ylabel style={font=\small},
    ymajorgrids=true,
    grid style={dashed,black!20},
    unbounded coords=discard,
    filter discard warning=false,
    legend style={at={(0.5,1.035)},anchor=south,legend columns=2,
                  font=\scriptsize,draw=none,fill=none},
    legend cell align=left,
    ymin=-5.5, ymax=2,
    ytick={-5,-4,...,2}
]

\addplot[black,thick,mark=*,mark size=1.5pt,forget plot,
         restrict x to domain=2019:2025]
    table[col sep=comma,x=year,y=balance_actualizacion]{datos/escenarios_4_1_4_2.csv};

% Inercial corresponde al escenario Actualización de la hoja 4.2.
\addplot[red,thick,dashed,mark=square*,mark size=1.5pt,restrict x to domain=2025:2035]
    table[col sep=comma,x=year,y=balance_actualizacion]{datos/escenarios_4_1_4_2.csv};
\addlegendentry{Inercial}

% Hasta 2027, la consolidación coincide con el escenario inercial.
\addplot[gray!80!black,thick,dotted,mark=triangle*,mark size=2pt,
         restrict x to domain=2027:2035]
    table[col sep=comma,x=year,y=balance_consolidacion]{datos/escenarios_4_1_4_2.csv};
\addlegendentry{Consolidación}

\addplot[black!45,thin,forget plot] coordinates {(2019,0) (2035,0)};
\node[anchor=west,font=\footnotesize,text=gray!80!black] at (axis cs:2035.5,1.25423454585) {1,3};
\node[anchor=west,font=\footnotesize,text=red] at (axis cs:2035.5,-3.3) {-3,3};
\end{axis}
\end{tikzpicture}

\par\vspace{0.05cm}
{\scriptsize Fuente: MHCP y cálculos propios.}\par
{\scriptsize Datos hasta 2025 en línea continua. Ambos escenarios coinciden hasta 2027.}
\note{[Sources] Figuras.xlsx, hoja 4.2, B2:R2 (años 2019 a 2035), B4:R4 (Actualización) y B5:R5 (Consolidación). Inercial corresponde a Actualización en el Excel. MHCP: balance primario del GNC; cálculos propios para los escenarios inercial y de consolidación. Los datos hasta 2025 son comunes a los escenarios. Las proyecciones comienzan en 2026 y los escenarios coinciden hasta 2027. La consolidación comienza en 2028. En 2035: inercial, -3,3\% del PIB; consolidación, 1,2542\%. Desde 2028, el balance primario de consolidación es 0,2 + 0,1 por la diferencia entre la deuda del año anterior y 55, en puntos del PIB.}
\end{frame}
